% Generated by the blueprint tool (blueprint-pdf/1) from a blueprint snapshot; edit the
% Markdown sources, not this file.  Compile with `latexmk -pdf` or by running
% `pdflatex` twice (the second run settles the contents and the page references).
% Images, if any, are read from assets/ next to this file.
\documentclass[11pt]{article}
\usepackage[T1]{fontenc}
\usepackage[utf8]{inputenc}
\usepackage{lmodern}
\usepackage{textcomp}
\usepackage[a4paper,margin=2.5cm,headheight=14pt]{geometry}
\usepackage{amsmath,amssymb,amsthm}
\usepackage{amscd}  % KaTeX's CD diagrams; part of amsmath, so always there
\usepackage{mathtools}
\usepackage{bm}
\usepackage{graphicx}
\usepackage{xcolor}
\usepackage{array,booktabs,tabularx}
\usepackage{fancyvrb}
\usepackage[normalem]{ulem}
\usepackage{microtype}
% Optional, and only with their fonts: a small TeX installation has the
% packages without the fonts, which fails at the first formula.
\IfFileExists{mathrsfs.sty}{\IfFileExists{rsfs10.tfm}{\usepackage{mathrsfs}}{}}{}
\IfFileExists{stmaryrd.sty}{\IfFileExists{stmary10.tfm}{\usepackage{stmaryrd}}{}}{}
\IfFileExists{cancel.sty}{\usepackage{cancel}}{}
\usepackage[hyphens]{url}
\usepackage[pdfusetitle,bookmarksnumbered=false]{hyperref}
\usepackage{bookmark}
\definecolor{bplink}{HTML}{1F4E9C}
\definecolor{bpmuted}{HTML}{5F676E}
\hypersetup{colorlinks=true,linkcolor=bplink,urlcolor=bplink,citecolor=bplink}
\setlength{\parindent}{0pt}
\setlength{\parskip}{0.55em plus 0.2em minus 0.1em}
% Lists as tight as the page's.
\makeatletter
\def\@listi{\leftmargin\leftmargini\parsep 0pt\topsep 0.3em plus 0.1em\itemsep 0.1em plus 0.1em}
\let\@listI\@listi
\def\@listii{\leftmargin\leftmarginii\labelwidth\leftmarginii\advance\labelwidth-\labelsep\parsep 0pt\topsep 0.1em\itemsep 0.1em}
% An ordered list that does not start at 1.
\newcommand{\bpstartat}[1]{\setcounter{\@enumctr}{\numexpr#1-1\relax}}
\makeatother
\setcounter{tocdepth}{3}
\emergencystretch=3em
\allowdisplaybreaks

% A numbered statement, set as a paper sets one: "Theorem 2.1 (Title)."
% #1 is the lead ("Theorem~2.1"), #2 what follows it in roman.
\newenvironment{bpstatement}[2]{\par\addvspace{0.6\baselineskip}\noindent\textbf{#1}#2\textbf{.}\enspace\ignorespaces}
  {\par\addvspace{0.4\baselineskip}}
\newenvironment{bpstatementit}[2]{\begin{bpstatement}{#1}{#2}\itshape\ignorespaces}{\end{bpstatement}}
% A step: a relation with prose of its own, read under the object it starts from.
\newenvironment{bpstep}[1]{\par\addvspace{0.3\baselineskip}\begin{list}{}{\leftmargin=1.6em\rightmargin=0pt\topsep=0pt}\item\relax
  {\small\color{bpmuted}#1}\par\nopagebreak}{\end{list}\par}
% Notes about where an object is written out.
\newenvironment{bpnote}{\par\small\color{bpmuted}\itshape}{\par}
\newcommand{\bpomitted}[1]{\par{\small\color{bpmuted}\itshape #1\par}}
\newcommand{\bpbodyheading}[1]{\par\addvspace{0.4\baselineskip}\noindent\textbf{#1}\par\nopagebreak}
\newcommand{\bprule}{\par\medskip\noindent\hfil\rule{0.3\linewidth}{0.4pt}\hfil\par\medskip}
\newcommand{\bpboundary}[1]{\par\nopagebreak{\small\color{bpmuted}#1\par}\nopagebreak}
% A reference to something this export leaves out, and one that resolves nowhere.
\newcommand{\bpout}[1]{#1\textsuperscript{\hyperlink{bp-outside}{\ensuremath{\ast}}}}
\newcommand{\bpbroken}[1]{\textcolor{red!70!black}{[#1?]}}
\newcommand{\bpslug}[1]{\texttt{#1}}
% A bibliography entry: label, anchor, entry; the label hangs in the margin.
\newcommand{\bpbibitem}[3]{\par\setbox0\hbox{[#1]\enspace}\ifdim\wd0<5em\setbox0\hbox to 5em{[#1]\hfil}\fi
  \hangindent=\wd0 \hangafter=1 \noindent\box0\phantomsection\label{#2}#3\par\smallskip}
% A character the export knows no way to set.
\newcommand{\bpglyph}[1]{\fbox{\scriptsize #1}}
% Images, at their own size unless they are wider than the line.
\newsavebox{\bpimagebox}
\newcommand{\bpimage}[1]{\sbox{\bpimagebox}{\includegraphics{#1}}%
  \ifdim\wd\bpimagebox>\linewidth\resizebox{\linewidth}{!}{\usebox{\bpimagebox}}\else\usebox{\bpimagebox}\fi}
\newcommand{\bpmissingimage}[1]{\fbox{\parbox{0.8\linewidth}{\centering\small\color{bpmuted}Image not included: #1}}}
\DefineVerbatimEnvironment{bpverbatim}{Verbatim}{fontsize=\small,xleftmargin=1em}
\renewcommand{\qedsymbol}{\ensuremath{\blacksquare}}
\hypersetup{pdftitle={Weil Conjectures — 17 The Deligne–Laumon theorem for curves},pdfsubject={Chapter 17 The Deligne–Laumon theorem for curves; all levels},pdfcreator={blueprint}}
% KaTeX commands the website knows
\providecommand{\R}{\mathbb{R}}\providecommand{\N}{\mathbb{N}}\providecommand{\Z}{\mathbb{Z}}
\providecommand{\Q}{\mathbb{Q}}\providecommand{\C}{\mathbb{C}}
\providecommand{\Reals}{\mathbb{R}}\providecommand{\reals}{\mathbb{R}}\providecommand{\Complex}{\mathbb{C}}
\providecommand{\cnums}{\mathbb{C}}\providecommand{\natnums}{\mathbb{N}}\providecommand{\Bbb}[1]{\mathbb{#1}}
\providecommand{\bold}[1]{\mathbf{#1}}\providecommand{\frak}[1]{\mathfrak{#1}}\providecommand{\mathbbm}[1]{\mathbb{#1}}
\providecommand{\lang}{\langle}\providecommand{\rang}{\rangle}\providecommand{\bra}[1]{\langle #1\rvert}
\providecommand{\ket}[1]{\lvert #1\rangle}\providecommand{\braket}[1]{\langle #1\rangle}
\providecommand{\set}[1]{\{#1\}}\providecommand{\Set}[1]{\left\{#1\right\}}
\providecommand{\infin}{\infty}\providecommand{\isin}{\in}\providecommand{\plusmn}{\pm}\providecommand{\sub}{\subset}
\providecommand{\sube}{\subseteq}\providecommand{\supe}{\supseteq}\providecommand{\exist}{\exists}
\providecommand{\bull}{\bullet}\providecommand{\larr}{\leftarrow}\providecommand{\rarr}{\rightarrow}
\providecommand{\lrarr}{\leftrightarrow}\providecommand{\Larr}{\Leftarrow}\providecommand{\Rarr}{\Rightarrow}
\providecommand{\Lrarr}{\Leftrightarrow}\providecommand{\uarr}{\uparrow}\providecommand{\darr}{\downarrow}
\providecommand{\harr}{\leftrightarrow}\providecommand{\lArr}{\Leftarrow}\providecommand{\rArr}{\Rightarrow}
\providecommand{\hArr}{\Leftrightarrow}\providecommand{\alef}{\aleph}\providecommand{\alefsym}{\aleph}
\providecommand{\real}{\Re}\providecommand{\image}{\Im}\providecommand{\weierp}{\wp}
\providecommand{\Alpha}{\mathrm{A}}\providecommand{\Beta}{\mathrm{B}}\providecommand{\Epsilon}{\mathrm{E}}
\providecommand{\Zeta}{\mathrm{Z}}\providecommand{\Eta}{\mathrm{H}}\providecommand{\Iota}{\mathrm{I}}
\providecommand{\Kappa}{\mathrm{K}}\providecommand{\Mu}{\mathrm{M}}\providecommand{\Nu}{\mathrm{N}}
\providecommand{\Omicron}{\mathrm{O}}\providecommand{\Rho}{\mathrm{P}}\providecommand{\Tau}{\mathrm{T}}
\providecommand{\Chi}{\mathrm{X}}\providecommand{\omicron}{o}\providecommand{\gt}{>}\providecommand{\lt}{<}
\providecommand{\And}{\&}\providecommand{\hdashline}{\hline}\providecommand{\operatornamewithlimits}[1]{\operatorname*{#1}}
\providecommand{\llbracket}{[\![}\providecommand{\rrbracket}{]\!]}\providecommand{\cancel}[1]{#1}
\providecommand{\mathscr}{\mathcal}
% the characters beyond ASCII this document uses
\DeclareUnicodeCharacter{2160}{\textup{I}}
\DeclareUnicodeCharacter{2161}{\textup{II}}
\DeclareUnicodeCharacter{2162}{\textup{III}}
% the project's macros ([katex.macros] in blueprint.toml)
\providecommand{\A}{}\renewcommand{\A}{\mathbf A}
\providecommand{\AEn}{}\renewcommand{\AEn}{A_{E,n}}
\providecommand{\Aaf}{}\renewcommand{\Aaf}{\mathbf A}
\providecommand{\Ab}{}\renewcommand{\Ab}{\mathsf{Ab}}
\providecommand{\AffSch}{}\renewcommand{\AffSch}{\operatorname{AffSch}}
\providecommand{\Aone}{}\renewcommand{\Aone}{\mathbf A^1}
\providecommand{\Bl}{}\renewcommand{\Bl}{\operatorname{Bl}}
\providecommand{\Br}{}\renewcommand{\Br}{\mathrm{Br}}
\providecommand{\C}{}\renewcommand{\C}{\mathbb C}
\providecommand{\Cech}{}\renewcommand{\Cech}{\check{H}}
\providecommand{\Coext}{}\renewcommand{\Coext}{\operatorname{Coext}}
\providecommand{\Cone}{}\renewcommand{\Cone}{\operatorname{Cone}}
\providecommand{\Cons}{}\renewcommand{\Cons}{\operatorname{Cons}}
\providecommand{\Cref}{}\renewcommand{\Cref}[1]{\text{#1}}
\providecommand{\D}{}\renewcommand{\D}{D}
\providecommand{\DLAFT}{}\renewcommand{\DLAFT}{\operatorname{FT}}
\providecommand{\DLAFr}{}\renewcommand{\DLAFr}{\operatorname{Fr}}
\providecommand{\DLAQ}{}\renewcommand{\DLAQ}{{\overline{\mathbf Q}_\ell}}
\providecommand{\DLRQ}{}\renewcommand{\DLRQ}{{\overline{\mathbf Q}_\ell}}
\providecommand{\Dbc}{}\renewcommand{\Dbc}{D^b_c}
\providecommand{\Dc}{}\renewcommand{\Dc}{D_{\mathrm{cons}}}
\providecommand{\Dctf}{}\renewcommand{\Dctf}{D^b_{ctf}}
\providecommand{\Dic}{}\renewcommand{\Dic}{D_{\IndCons}}
\providecommand{\Dplus}{}\renewcommand{\Dplus}{D^{+}}
\providecommand{\End}{}\renewcommand{\End}{\operatorname{End}}
\providecommand{\Et}{}\renewcommand{\Et}{\mathrm{\acute{e}t}}
\providecommand{\Ex}{}\renewcommand{\Ex}{\mathrm{Ex}}
\providecommand{\Ext}{}\renewcommand{\Ext}{\operatorname{Ext}}
\providecommand{\F}{}\renewcommand{\F}{\mathbb F}
\providecommand{\FT}{}\renewcommand{\FT}{\operatorname{FT}}
\providecommand{\Fbar}{}\renewcommand{\Fbar}{{\overline{\mathbf F}_q}}
\providecommand{\Fell}{}\renewcommand{\Fell}{\mathbf F_\ell}
\providecommand{\Fib}{}\renewcommand{\Fib}{\operatorname{Fib}}
\providecommand{\Fix}{}\renewcommand{\Fix}{\operatorname{Fix}}
\providecommand{\Fp}{}\renewcommand{\Fp}{{\mathbf F_{p}}}
\providecommand{\Fq}{}\renewcommand{\Fq}{{\mathbf F_q}}
\providecommand{\Fqn}{}\renewcommand{\Fqn}{{\mathbf F_{q^{n}}}}
\providecommand{\Frob}{}\renewcommand{\Frob}{\operatorname{Frob}}
\providecommand{\GL}{}\renewcommand{\GL}{\operatorname{GL}}
\providecommand{\Gal}{}\renewcommand{\Gal}{\operatorname{Gal}}
\providecommand{\Hens}{}\renewcommand{\Hens}{\operatorname{Hens}}
\providecommand{\Hilb}{}\renewcommand{\Hilb}{\operatorname{Hilb}}
\providecommand{\Hom}{}\renewcommand{\Hom}{\operatorname{Hom}}
\providecommand{\IndCons}{}\renewcommand{\IndCons}{\operatorname{IndCons}}
\providecommand{\Kar}{}\renewcommand{\Kar}{\operatorname{Kar}}
\providecommand{\LT}{}\renewcommand{\LT}{\operatorname{LT}}
\providecommand{\Lie}{}\renewcommand{\Lie}{\operatorname{Lie}}
\providecommand{\Loc}{}\renewcommand{\Loc}{\operatorname{Loc}}
\providecommand{\Mod}{}\renewcommand{\Mod}{\mathsf{Mod}}
\providecommand{\N}{}\renewcommand{\N}{\mathbb{N}}
\providecommand{\Needspace}{}\renewcommand{\Needspace}{\par}
\providecommand{\O}{}\renewcommand{\O}{{\mathcal{O}}}
\providecommand{\Ob}{}\renewcommand{\Ob}{\mathrm{Ob}}
\providecommand{\Oe}{}\renewcommand{\Oe}{\mathcal O_E}
\providecommand{\PP}{}\renewcommand{\PP}{\mathbf{P}}
\providecommand{\Perf}{}\renewcommand{\Perf}{\operatorname{Perf}}
\providecommand{\Pic}{}\renewcommand{\Pic}{\mathrm{Pic}}
\providecommand{\Pm}{}\renewcommand{\Pm}{\mathbf P^{m}}
\providecommand{\Pone}{}\renewcommand{\Pone}{\mathbf P^{1}}
\providecommand{\Pqn}{}\renewcommand{\Pqn}{\mathbf P^{n}}
\providecommand{\ProFin}{}\renewcommand{\ProFin}{\operatorname{ProFin}}
\providecommand{\Proet}{}\renewcommand{\Proet}{\mathrm{pro\acute{e}t}}
\providecommand{\Q}{}\renewcommand{\Q}{\mathbb Q}
\providecommand{\Qb}{}\renewcommand{\Qb}{\overline{\mathbf Q}_{\ell}}
\providecommand{\Qbar}{}\renewcommand{\Qbar}{\overline{\mathbf Q}}
\providecommand{\QbarX}{}\renewcommand{\QbarX}{\overline{\mathbf Q}_{\ell,X}}
\providecommand{\Qbarl}{}\renewcommand{\Qbarl}{{\overline{\mathbf Q}_{\ell}}}
\providecommand{\Qell}{}\renewcommand{\Qell}{{\mathbf{Q}_{\ell}}}
\providecommand{\Ql}{}\renewcommand{\Ql}{\mathbf Q_{\ell}}
\providecommand{\Qlb}{}\renewcommand{\Qlb}{{\overline{\mathbf Q}_\ell}}
\providecommand{\Qlbar}{}\renewcommand{\Qlbar}{\overline{\mathbf Q}_\ell}
\providecommand{\R}{}\renewcommand{\R}{\mathbb{R}}
\providecommand{\RG}{}\renewcommand{\RG}{\operatorname{\mathrm{R}\Gamma}}
\providecommand{\RGamma}{}\renewcommand{\RGamma}{R\Gamma}
\providecommand{\RGammac}{}\renewcommand{\RGammac}{R\Gamma_c}
\providecommand{\RHom}{}\renewcommand{\RHom}{R\mathcal{H}om}
\providecommand{\RHomX}{}\renewcommand{\RHomX}{R\mathcal H om_X}
\providecommand{\Rat}{}\renewcommand{\Rat}{\operatorname{IsRat}}
\providecommand{\Res}{}\renewcommand{\Res}{\operatorname{Res}}
\providecommand{\Rlim}{}\renewcommand{\Rlim}{R\!\varprojlim}
\providecommand{\SIm}{}\renewcommand{\SIm}{\operatorname{SIm}}
\providecommand{\Sch}{}\renewcommand{\Sch}{\mathsf{Sch}}
\providecommand{\Sh}{}\renewcommand{\Sh}{\operatorname{Sh}}
\providecommand{\Span}{}\renewcommand{\Span}{\operatorname{span}}
\providecommand{\Spec}{}\renewcommand{\Spec}{\operatorname{Spec}}
\providecommand{\Tr}{}\renewcommand{\Tr}{\operatorname{Tr}}
\providecommand{\Trsm}{}\renewcommand{\Trsm}{\operatorname{Tr}^{\mathrm{sm}}}
\providecommand{\Wt}{}\renewcommand{\Wt}{\operatorname{Wt}}
\providecommand{\Z}{}\renewcommand{\Z}{\mathbb{Z}}
\providecommand{\Zar}{}\renewcommand{\Zar}{\mathrm{zar}}
\providecommand{\Zell}{}\renewcommand{\Zell}{{\mathbf{Z}_{\ell}}}
\providecommand{\Zl}{}\renewcommand{\Zl}{{\mathbf Z_\ell}}
\providecommand{\affet}{}\renewcommand{\affet}[1]{{#1}^{\mathrm{aff}}_{\mathrm{ét}}}
\providecommand{\affproet}{}\renewcommand{\affproet}[1]{{#1}^{\mathrm{aff}}_{\mathrm{proét}}}
\providecommand{\boxt}{}\renewcommand{\boxt}{\boxtimes_A^{\mathrm L}}
\providecommand{\cF}{}\renewcommand{\cF}{\mathcal{F}}
\providecommand{\cHom}{}\renewcommand{\cHom}{R\mathcal{H}\!om}
\providecommand{\cI}{}\renewcommand{\cI}{\mathcal{I}}
\providecommand{\cL}{}\renewcommand{\cL}{\mathcal L}
\providecommand{\cM}{}\renewcommand{\cM}{\mathcal M}
\providecommand{\cN}{}\renewcommand{\cN}{\mathcal N}
\providecommand{\cO}{}\renewcommand{\cO}{\mathcal O}
\providecommand{\cT}{}\renewcommand{\cT}{\mathcal T}
\providecommand{\checked}{}\renewcommand{\checked}{\color{red}\aftergroup\restorecolor}
\providecommand{\coeff}{}\renewcommand{\coeff}{\operatorname{coeff}}
\providecommand{\coker}{}\renewcommand{\coker}{\operatorname{coker}}
\providecommand{\colim}{}\renewcommand{\colim}{\operatorname{colim}}
\providecommand{\constructedby}{}\renewcommand{\constructedby}[1]{\par\smallskip\noindent\textcolor{purple!60!black}{\textsf{Constructed by:} }
  \footnotesize\sloppy\nolinkurl{#1}\par}
\providecommand{\deepresult}{}\renewcommand{\deepresult}[1]{\par\smallskip\noindent\textcolor{orange!45!black}{
  \textsf{Blueprint role: named deep theorem node; see the cluster proof architecture.}}
  \par\noindent\textcolor{orange!45!black}{\textsf{Mathematical source:} }#1\par}
\providecommand{\dt}{}\renewcommand{\dt}{\mathbin{\otimes^{\mathbf L}}}
\providecommand{\dw}{}\renewcommand{\dw}{\operatorname{dw}}
\providecommand{\et}{}\renewcommand{\et}[1]{{#1}_{\mathrm{ét}}}
\providecommand{\ev}{}\renewcommand{\ev}{\operatorname{ev}}
\providecommand{\fib}{}\renewcommand{\fib}{\operatorname{fib}}
\providecommand{\formalizationnote}{}\renewcommand{\formalizationnote}[1]{\begin{quote}\small\raggedright\textbf{Formalization note.} #1\end{quote}}
\providecommand{\fpqc}{}\renewcommand{\fpqc}{X_{\mathrm{fpqc}}}
\providecommand{\gammacont}{}\renewcommand{\gammacont}{\Gamma_{\mathrm{cont}}}
\providecommand{\gm}{}\renewcommand{\gm}{\mathbf G_m}
\providecommand{\hocolim}{}\renewcommand{\hocolim}{\operatorname{hocolim}}
\providecommand{\holim}{}\renewcommand{\holim}{\operatorname{holim}}
\providecommand{\hrB}{}\renewcommand{\hrB}{\@ifnextchar]{\hrEnd}{\hrStep}}
\providecommand{\hrEnd}{}\renewcommand{\hrEnd}{\@firstoftwo{}}
\providecommand{\hrStep}{}\renewcommand{\hrStep}{\@firstoftwo{\hrB}}
\providecommand{\id}{}\renewcommand{\id}{\operatorname{id}}
\providecommand{\kunFl}{}\renewcommand{\kunFl}{\mathbf F_\ell}
\providecommand{\kunZn}{}\renewcommand{\kunZn}{\mathbf Z/n\mathbf Z}
\providecommand{\kuncd}{}\renewcommand{\kuncd}{\operatorname{cd}}
\providecommand{\kundim}{}\renewcommand{\kundim}{\dim}
\providecommand{\kunotimes}{}\renewcommand{\kunotimes}[1]{\otimes^{\mathbf L}_{#1}}
\providecommand{\kunsupp}{}\renewcommand{\kunsupp}{\operatorname{Supp}}
\providecommand{\lotimes}{}\renewcommand{\lotimes}{\otimes^{\mathrm L}}
\providecommand{\m}{}\renewcommand{\m}{{\mathfrak{m}}}
\providecommand{\n}{}\renewcommand{\n}{{\mathfrak{n}}}
\providecommand{\nagataGap}{}\renewcommand{\nagataGap}[1]{\par\noindent\textbf{Formalization gap.} #1\par}
\providecommand{\openobligation}{}\renewcommand{\openobligation}[1]{\par\smallskip\noindent\textcolor{red!65!black}{\textsf{Open obligation:} }
  \footnotesize\sloppy\nolinkurl{#1}\par}
\providecommand{\p}{}\renewcommand{\p}{{\mathfrak{p}}}
\providecommand{\pr}{}\renewcommand{\pr}{\operatorname{pr}}
\providecommand{\preshv}{}\renewcommand{\preshv}[1]{\mathrm{PreShv}(#1)}
\providecommand{\pro}{}\renewcommand{\pro}[1]{\mathrm{Pro}(#1)}
\providecommand{\proet}{}\renewcommand{\proet}{X_{\mathrm{proét}}}
\providecommand{\q}{}\renewcommand{\q}{{\mathfrak{q}}}
\providecommand{\rad}{}\renewcommand{\rad}{\operatorname{rad}}
\providecommand{\rank}{}\renewcommand{\rank}{\operatorname{rank}}
\providecommand{\red}{}\renewcommand{\red}{\color{red}}
\providecommand{\restorecolor}{}\renewcommand{\restorecolor}{\color{black}}
\providecommand{\retiredobligation}{}\renewcommand{\retiredobligation}[1]{\par\smallskip\noindent\textcolor{red!75!black}{\textsf{Retired invalid route:} }
  \footnotesize\sloppy\nolinkurl{#1}\par}
\providecommand{\rk}{}\renewcommand{\rk}{\operatorname{rk}}
\providecommand{\sep}{}\renewcommand{\sep}{\textnormal{sep}}
\providecommand{\shv}{}\renewcommand{\shv}[1]{\mathrm{Shv}(#1)}
\providecommand{\spec}{}\renewcommand{\spec}[1]{\mathrm{Spec}(#1)}
\providecommand{\stacks}{}\renewcommand{\stacks}[1]{\href{https://stacks.math.columbia.edu/tag/#1}{\textnormal{Stacks, Tag~#1}}}
\providecommand{\stacksproject}{}\renewcommand{\stacksproject}[1]{{\cite[\href{https://stacks.math.columbia.edu/tag/#1}{Tag #1}]{stacks-project}}}
\providecommand{\statusnote}{}\renewcommand{\statusnote}[1]{\begin{quote}\small\raggedright\textbf{Status.} #1\end{quote}}
\providecommand{\tildeideal}{}\renewcommand{\tildeideal}[2]{\widetilde{(#1, #2)}}
\providecommand{\tildering}{}\renewcommand{\tildering}[3]{{#1}_{\widetilde{(#2, #3)}}}
\providecommand{\tr}{}\renewcommand{\tr}{\operatorname{tr}}
\providecommand{\wt}{}\renewcommand{\wt}{\operatorname{wt}}

\begin{document}
\pagestyle{plain}
\begin{center}
{\LARGE\bfseries Weil Conjectures\par}
\vspace{0.7em}
{\large Chapter 17: The Deligne–Laumon theorem for curves\par}
\vspace{0.5em}
{\small\color{bpmuted}Snapshot 2026-09-26, revision 8e4cb88e2a\par}
\end{center}
\vspace{0.5em}
{\small\tableofcontents}
\bigskip\hrule\bigskip

% object: chap-The-Deligne-Laumon-theorem-for-curves (blueprint/chap-The-Deligne-Laumon-theorem-for-curves/_section.md)
\section*{\texorpdfstring{17\quad The Deligne–Laumon theorem for curves}{17 The Deligne–Laumon theorem for curves}}
\phantomsection\label{bp1}%
\addcontentsline{toc}{section}{\texorpdfstring{17\quad The Deligne–Laumon theorem for curves}{17 The Deligne–Laumon theorem for curves}}
% text: prose-17-01
\phantomsection\label{bp2}%
Let \(F\) be a pure lisse Weil sheaf of weight \(w\) on a smooth curve \(C\). Why should \(H^1_c(C,F)\) have weights at most \(w+1\)? For the constant sheaf on a smooth projective curve this is Weil's theorem, proved with the Jacobian, and for an arbitrary \(F\) there is no such object to work with. Yet the Deligne–Laumon theorem comes down to exactly this bound. The theorem says that \(H^r_c(C,F)\) has weights at most \(w+r\) (Theorem \hyperref[bp47]{17.15}), and only \(r=1\) requires work. Indeed, on an affine curve \(H^0_c\) vanishes, while by Lemma \bpout{15.7} \(H^2_c\) is a twist of the coinvariants of geometric monodromy, pure of weight \(w+2\). For \(r=1\) the bound can be strict: on \(\mathbf G_m\) the group \(H^1_c\) comes from the boundary and has weight \(0\) (see the example of Chapter \bpout{13}).

For a Gauss sum the answer takes two steps. The transform \(\hat\chi(y)=\overline\chi(y)g\) has the same absolute value for all \(y\ne0\), and Plancherel finds the common value \(\sqrt q\). On an open of the affine line, Chapter \bpout{16} makes the same two steps available for sheaves. More precisely, let \(M\) be irreducible on a dense open \(j:U_0\hookrightarrow\mathbf A^1\), with no geometric constituent of the form \(\mathcal L_{\psi(ax)}\), and let \(M'\) be the lisse sheaf with \(\operatorname{FT}_\psi((j_*M)[1])\simeq(j'_*M')[1]\) (Proposition \bpout{16.17}). Then

\[
H^1_c(\mathbf A^1,j_*M)\simeq(j'_*M')_{0'}
\]

(Corollary \hyperref[bp19]{17.6}). Since a stalk of a middle extension has weights at most the weight of the lisse sheaf (Proposition \bpout{15.5}), it suffices to show that \(M'\) is pure of weight \(w+1\). Here the second step is the easy one, because once \(M'\) is pure, the sums of squared traces of Chapter \bpout{15} give its weight \(w+1\) (Theorem \hyperref[bp16]{17.5}).

The first step, purity of \(M'\), is the deep one. The only purity theorem we have is Theorem \bpout{14.26}, for real sheaves. But the transform of a real sheaf is not real, since \(\psi\) takes nonreal values. For example, over \(\mathbf F_3\) let \(M\) be the Kummer sheaf of the quadratic character \(\chi\) on \(\mathbf G_m\), which is real. By the computation of Chapter \bpout{16} and the signs in (16.3), the trace of \(M'\) at \(y\in\mathbf F_3^\times\) is \(-\chi(y)g\). Here \(g=\sum_x\chi(x)\psi(x)=\omega-\omega^2=1+2\omega\) is a square root of \(-3\), which is not real. In this case, admittedly, a twist by \(\mathcal B_{1/g}\) makes \(M'\) real, but in general no twist is known to do this.

The plan is to realize \(M'\), up to a twist, as a quotient of a real sheaf on the dual line, to which that theorem applies. To build it we add a conjugate. First we twist \(M\) to weight zero, \(M^0=M\otimes\mathcal B_b\). Then we transform \(j_!M^0\) with the kernel \(\mathcal K_\psi\) and \(j_!(M^0)^\vee\) with the kernel \(\mathcal K_{\psi^{-1}}\). At each point the traces of the two transforms are complex conjugate, so their sum is a real sheaf \(H\) on the dual line. Since the twist \(M'\otimes\mathcal B_b\) is a quotient of \(H\) (Proposition \hyperref[bp12]{17.4}), purity of real sheaves then makes \(M'\) pure. We use \(j_!\) rather than \(j_*\) because its transform is a sheaf, and by Lemma \hyperref[bp5]{17.1} this transform differs from the transform of \(j_*M\) by a lisse sheaf of weights at most \(w\).

Passing from one irreducible \(M\) to every pure sheaf needs care, because \(j_*\) is not exact in general. For example, the extension \(0\to\overline{\mathbf Q}_\ell(1)\to F\to\overline{\mathbf Q}_\ell\to0\) on \(\mathbf G_m\) given by the Kummer class of the coordinate is lisse, but \(j_*F\to j_*\overline{\mathbf Q}_\ell\) is not surjective at \(0\), since the monodromy of \(F\) around \(0\) is a nontrivial unipotent matrix. For sheaves of one weight \(w\), however, exactness holds by a weight gap (Proposition \hyperref[bp27]{17.9}): by local duality at the boundary points (Lemma \hyperref[bp23]{17.7}), the connecting map \(j_*F_2\to R^1j_*F_1\) goes from weights at most \(w\) to weights at least \(w+2\), and so it vanishes.

% object: sec-Purity-of-the-lisse-part (blueprint/chap-The-Deligne-Laumon-theorem-for-curves/sec-Purity-of-the-lisse-part/_section.md), content hidden
\subsection*{\texorpdfstring{17.1\quad Purity of the lisse part}{17.1 Purity of the lisse part}}
\phantomsection\label{bp3}%
\addcontentsline{toc}{subsection}{\texorpdfstring{17.1\quad Purity of the lisse part}{17.1 Purity of the lisse part}}
% text: prose-17-02
\phantomsection\label{bp4}%
Throughout this section and the next, \(j:U_0\hookrightarrow\mathbf A^1_{k_0}\) is a dense open and \(M\) is a nonzero irreducible lisse Weil sheaf on \(U_0\), pure of weight \(w\). We assume that no geometric irreducible constituent of \(M\) is the restriction of \(\mathcal L_{\psi(ax)}\) for some \(a\in\bar k_0\), and we let \(M'\) be the lisse Weil sheaf on a dense open \(j':U'_0\hookrightarrow\mathbf A^{1\vee}_{k_0}\) with \(\operatorname{FT}_\psi((j_*M)[1])\simeq(j'_*M')[1]\) (Proposition \bpout{16.17}).

We build the real sheaf from extensions by zero, because for them we can check reality on the trace function. Indeed, at a point of \(U_0\) the eigenvalues of \(M^0\) and of its dual have absolute value one, and hence they are complex conjugate. At a boundary point, however, the stalk of \(j_*M^0\) is the space of inertia invariants, and the stalk of \(j_*(M^0)^\vee\) is the space of invariants of the dual representation, but nothing we know at this stage makes their Frobenius eigenvalues conjugate. With \(j_!\), by contrast, both stalks vanish there, and the trace function of the sum of the two transforms is clearly real.

The price is that \(j_!M\) is not a simple perverse sheaf, so its transform \(G\) is not \(j'_*M'\) but an extension of \(j'_*M'\) by the transform \(B'\) of the boundary. This \(B'\) is a lisse sheaf of weights at most \(w\) (Lemma \hyperref[bp5]{17.1}), and in the proof of Proposition \hyperref[bp12]{17.4} we use only that \(G\) is a sheaf and that it surjects onto \(j'_*M'\).

% object: lemma-fourier-transform-boundary-sequence (blueprint/chap-The-Deligne-Laumon-theorem-for-curves/sec-Purity-of-the-lisse-part/lemma-fourier-transform-boundary-sequence/_section.md)
\phantomsection\label{bp5}%
\begin{bpstatementit}{Lemma~17.1}{ (Fourier transform of the boundary sequence)}
Put \(B=j_*M/j_!M\). There are sheaves \(G\) and \(B'\) on the dual line such that \(\operatorname{FT}_\psi(j_!M)=G\), \(\operatorname{FT}_\psi(B)=B'[1]\), and the sequence

\[
0\longrightarrow B'\longrightarrow G
\longrightarrow j'_*M'\longrightarrow0
\]

is exact. The sheaf \(B'\) is lisse on the entire dual line and has weights at most \(w\).
\end{bpstatementit}

% proof: proof/lemma-fourier-transform-boundary-sequence
\begin{proof}[{Proof}]\phantomsection\label{bp6}%
The three assertions are of different kinds. That the transform of \(j_!M\) is a single sheaf is a vanishing statement for the cohomology of the twisted sheaves \(M\otimes\mathcal L_{\psi(yx)}\) on \(U\), whereas the exact sequence comes from transforming the boundary sequence \(0\to j_!M\to j_*M\to B\to0\) term by term. Finally, the bound on \(B'\) comes from the stalks of \(B\), which are known from Proposition \bpout{15.5}.

For \(G\), the stalk formula of Proposition \bpout{16.5} says that the stalk of \(\operatorname{FT}_\psi(j_!M)\) at \(y\) is \(R\Gamma_c(U,M\otimes\mathcal L_{\psi(yx)})[1]\), so we must show that only \(H^1_c\) is nonzero. We apply Lemma \bpout{15.7} to the lisse sheaf \(M\otimes\mathcal L_{\psi(yx)}\) on the affine curve \(U\), with \(j\) the identity. It shows that the group in degree zero vanishes and that the group in degree two is a twist of the coinvariants of geometric monodromy on \(M\otimes\mathcal L_{\psi(yx)}\). But a nonzero coinvariant would be a geometric quotient of \(M\) isomorphic to \(\mathcal L_{\psi(-yx)}\), and our assumption on \(M\) excludes such a constituent, so the group in degree two vanishes as well. The shift then places the remaining group \(H^1_c\) in degree zero, and the transform is a sheaf \(G\).

The transform of \(B\) is computed point by point. After geometric base change, the punctual sheaf \(B\) is a finite sum of skyscrapers with coefficient modules, and the transform of each skyscraper is a linear character sheaf tensored with that module and shifted by one. Since a linear character sheaf is lisse on the whole dual line, \(\operatorname{FT}_\psi(B)=B'[1]\) with \(B'\) lisse.

Applying the Fourier transform to the triangle obtained from \(0\to j_!M\to j_*M\to B\to0\), and using Proposition \bpout{16.17} for the middle term, we obtain a triangle \(G\to j'_*M'\to B'[1]\to\). Here \(G\) and \(j'_*M'\) are sheaves in degree zero, and \(B'[1]\) is a sheaf in degree \(-1\). Hence the long exact sequence of cohomology sheaves has only three nonzero terms, and it is the stated exact sequence.

The weights of \(B'\) come from those of \(B\), because a character sheaf has weight zero, and the weights of \(B\) are at most \(w\) by Proposition \bpout{15.5}. Only at boundary points that are not rational does this need an argument. So let \(s\) be a boundary point of degree \(d\) and \(y\) a rational point of the dual line. The geometric points over \(s\) contribute \(d\) summands to \(B'_{\bar y}\), which \(\operatorname{Frob}_y\) permutes cyclically, and its \(d\)th power acts on one summand as \(\operatorname{Frob}_s\) on \(B_{\bar s}\) times a root of unity coming from the character sheaf. Every eigenvalue \(\beta\) of \(\operatorname{Frob}_y\) on these summands therefore satisfies \(\beta^d=\zeta\alpha\), with \(\zeta\) a root of unity and \(\alpha\) an eigenvalue of \(\operatorname{Frob}_s\). Since \(|\iota\alpha|\le q^{dw/2}\), we get \(|\iota\beta|\le q^{w/2}\). At a point \(y\) of higher degree we replace \(q\) by the cardinality of \(\kappa(y)\).

For a filtration of \(B'\) with pure quotients, choose for each boundary point, over its residue field, a Frobenius-stable full flag of the coefficient module, which exists because \(\overline{\mathbf Q}_\ell\) is algebraically closed. Each successive quotient is then a line, on which Frobenius acts by one of its eigenvalues. Tensoring with the linear character sheaves and summing over the Frobenius orbit, we obtain a finite Weil filtration of \(B'\) whose quotients are pure of the weights of the coefficient eigenvalues, and these weights are at most \(w\).
\end{proof}

% object: example-kummer-no-boundary (blueprint/chap-The-Deligne-Laumon-theorem-for-curves/sec-Purity-of-the-lisse-part/example-kummer-no-boundary/_section.md)
\phantomsection\label{bp7}%
\begin{bpstatement}{Example~17.2}{ (A Kummer sheaf has no boundary term)}
Let \(M\) be the Kummer sheaf of a nontrivial character \(\chi\) on \(U_0=\mathbf G_m\). It is irreducible and pure of weight \(0\). Since it is tame at infinity, it has no geometric constituent \(\mathcal L_{\psi(ax)}\), and so it satisfies the hypotheses of this section. Inertia at the origin acts on \(M\) through \(\chi\), nontrivially, and hence \((j_*M)_{\bar0}=0\) by Lemma \bpout{15.2}. This means that \(j_!M=j_*M\) and \(B=0\), and the lemma gives \(G\simeq j'_*M'\). On \(\mathbf F_q\)-points this is visible: by (16.3) the trace of \(G\) at \(y\ne0\) is \(-\sum_x\chi(x)\psi(xy)=-\overline\chi(y)g\), with \(g=\sum_x\chi(x)\psi(x)\), and for the quadratic character this is the trace of \(M'\) computed at the start of Chapter \hyperref[bp1]{17}. At \(y=0\), on the other hand, the trace of \(G\) is \(-\sum_x\chi(x)=0\), in agreement with \(H^1_c(\mathbf G_m,M)=0\).
\end{bpstatement}

% text: prose-17-03
\phantomsection\label{bp8}%
The argument of Proposition \bpout{15.9} also applies to twists by the Fourier kernel, and we need this in the proof of Proposition \hyperref[bp12]{17.4}.

% object: lemma-twisted-envelope (blueprint/chap-The-Deligne-Laumon-theorem-for-curves/sec-Purity-of-the-lisse-part/lemma-twisted-envelope/_section.md)
\phantomsection\label{bp9}%
\begin{bpstatementit}{Lemma~17.3}{ (Twisting and adding the dual)}
Let \(\mathcal A\) be lisse and pure of weight zero, and let \(\mathcal L\) be lisse of rank one with local eigenvalues of finite order. Then

\[
(\mathcal A\otimes\mathcal L)
\oplus(\mathcal A^\vee\otimes\mathcal L^{-1})
\]

is real and pure of weight zero, and its extension by zero along an open immersion is real as well.
\end{bpstatementit}

% proof: proof/lemma-twisted-envelope
\begin{proof}[{Proof}]\phantomsection\label{bp10}%
As in the proof of Proposition \bpout{15.9}, the local eigenvalues of \(\mathcal A\otimes\mathcal L\) have absolute value one after applying \(\iota\), and those of the second summand are their inverses, hence their complex conjugates. The extension by zero is real as well, because its stalks outside the open are zero, and so its local characteristic polynomials there are equal to \(1\).
\end{proof}

% text: prose-17-04
\phantomsection\label{bp11}%
In the next proof we use Lemma \hyperref[bp5]{17.1} twice. Applying its vanishing argument to \(M^0\) and to its dual, we see that the transform of the real sheaf is a single sheaf. Then, using the surjection \(G\to j'_*M'\) of the lemma, we exhibit a twist of \(M'\) as a quotient of that real sheaf.

% object: proposition-purity-lisse-part-fourier-transform (blueprint/chap-The-Deligne-Laumon-theorem-for-curves/sec-Purity-of-the-lisse-part/proposition-purity-lisse-part-fourier-transform/_section.md)
\phantomsection\label{bp12}%
\begin{bpstatementit}{Proposition~17.4}{ (Purity of the lisse part of the Fourier transform)}
The sheaf \(M'\) is pure of some real weight \(v\) on the whole open where it is lisse.
\end{bpstatementit}

% proof: proof/proposition-purity-lisse-part-fourier-transform
\begin{proof}[{Proof}]\phantomsection\label{bp13}%
We twist to weight zero and add the dual, as in Proposition \bpout{15.9}, before taking the transform, so that a twist of \(M'\) becomes a quotient of a real sheaf, to which Theorem \bpout{14.26} applies.

Choose \(b\) as in Proposition \bpout{15.9}, so that \(M^0=M\otimes\mathcal B_b\) is pure of weight zero.

On the product of the line and the dual line, consider the sheaf

\[
F=
(p_1^*j_!M^0\otimes\mathcal K_\psi)
\ \oplus\
(p_1^*j_!(M^0)^\vee\otimes\mathcal K_{\psi^{-1}}).
\]

At each closed point, the Frobenius eigenvalues of the second summand are the complex conjugates of those of the first, because the eigenvalues of \(M^0\) have absolute value one and those of the character sheaves are roots of unity, while at the boundary both zero extensions have zero stalks. Hence \(F\) is real, as in Lemma \hyperref[bp9]{17.3}. Moreover, since duality replaces \(a\) by \(-a\), the dual \((M^0)^\vee\) also has no geometric constituent of the form \(\mathcal L_{\psi(ax)}\), so the argument of Lemma \hyperref[bp5]{17.1} shows that the direct image with compact support of \(F\) along \(p_2\) is concentrated in degree one. We set \(H=R^1p_{2!}F\).

Let \(y\) be a point with values in a finite field. Since \(Rp_{2!}F=H[-1]\), the trace formula for \(Rp_{2!}\) (Corollary \bpout{12.26}) gives \(-t_H(y)=\sum_x t_F(x,y)\). For each \(x\) the traces of the two summands of \(F\) at \((x,y)\) are complex conjugate, so \(t_H(y)\) is real. Over the finite extensions of \(\kappa(y)\) the same argument shows that the trace of every power of the local Frobenius is real. These traces are the power sums of the local eigenvalues, and the Newton identities express the coefficients of the characteristic polynomial through the power sums with rational coefficients. Hence the local characteristic polynomials of \(H\) have real coefficients, that is, \(H\) is real.

The first summand of \(H\) is \(R^1p_{2!}(p_1^*j_!M\otimes\mathcal K_\psi)\otimes\mathcal B_b\) by the projection formula for the constant twist, that is, \(G\otimes\mathcal B_b\) with \(G=\operatorname{FT}_\psi(j_!M)\) as in Lemma \hyperref[bp5]{17.1}. Tensoring the surjection \(G\to j'_*M'\) of that lemma with \(\mathcal B_b\), we see that \(M'\otimes\mathcal B_b\) is a quotient of a direct summand of \(H\). To apply the purity theorem we restrict to a dense open on which both \(H\) and \(M'\) are lisse. There \(M'\otimes\mathcal B_b\) is a lisse quotient of the real lisse sheaf \(H\), and it is irreducible by Lemma \bpout{16.14}, so by Theorem \bpout{14.26} it is pure of some real weight. Removing the twist by \(\mathcal B_b\), we find that \(M'\) is pure of some weight \(v\) on that open. Since the open where \(M'\) is lisse is a smooth geometrically connected curve, the last assertion of Corollary \bpout{15.6} extends the purity to all of it.
\end{proof}


% object: sec-The-weight-of-the-Fourier-transform (blueprint/chap-The-Deligne-Laumon-theorem-for-curves/sec-The-weight-of-the-Fourier-transform/_section.md), content hidden
\subsection*{\texorpdfstring{17.2\quad The weight of the Fourier transform}{17.2 The weight of the Fourier transform}}
\phantomsection\label{bp14}%
\addcontentsline{toc}{subsection}{\texorpdfstring{17.2\quad The weight of the Fourier transform}{17.2 The weight of the Fourier transform}}
% text: prose-17-05
\phantomsection\label{bp15}%
By Proposition \hyperref[bp12]{17.4}, \(M'\) is pure of some weight \(v\), and Plancherel now finds \(v\). Purity had to come first, since the growth exponent \(\nu\) of Definition \bpout{15.10} is the weight only for pure sheaves.

% object: theorem-weight-fourier-transform (blueprint/chap-The-Deligne-Laumon-theorem-for-curves/sec-The-weight-of-the-Fourier-transform/theorem-weight-fourier-transform/_section.md)
\phantomsection\label{bp16}%
\begin{bpstatementit}{Theorem~17.5}{ (The weight of the Fourier transform)}
The sheaf \(M'\) is pure of weight exactly \(w+1\).
\end{bpstatementit}

% proof: proof/theorem-weight-fourier-transform
\begin{proof}[{Proof}]\phantomsection\label{bp17}%
Let \(v\) be the weight of \(M'\) given by Proposition \hyperref[bp12]{17.4}. Shifts do not change the squared absolute values of traces, and, if a nonzero lisse sheaf on a dense open of the line is pure of weight \(u\), then its direct image \(j_*\) has \(\nu=u\) (Corollary \bpout{15.15}). Applying this to \(M'\) and to \(M\), and using the Plancherel formula (Corollary \bpout{16.8}), we get

\[
v=\nu(j'_*M')
=\nu(\operatorname{FT}_\psi((j_*M)[1]))
=\nu((j_*M)[1])+1=w+1.
\]

The first equality holds because \(M'\) is pure of weight \(v\), the second is the defining isomorphism of \(M'\), the third is Plancherel, and the last holds because \(M\) is pure of weight \(w\).
\end{proof}

% text: prose-17-06
\phantomsection\label{bp18}%
In the example of the chapter opening, \(M\) is the Kummer sheaf of the quadratic character over \(\mathbf F_3\), of weight \(0\), and the trace of \(M'\) at \(y\in\mathbf F_3^\times\) is \(-\chi(y)g\) with \(g^2=-3\). Its absolute value is \(\sqrt3=3^{1/2}\), which is weight \(1=w+1\), as the theorem predicts.

The weight of \(M'\) bounds \(H^1_c\), because the stalk of the transform at the origin of the dual line is the compactly supported cohomology of \(j_*M\).

% object: corollary-h-1-c-stalk-fourier (blueprint/chap-The-Deligne-Laumon-theorem-for-curves/sec-The-weight-of-the-Fourier-transform/corollary-h-1-c-stalk-fourier/_section.md)
\phantomsection\label{bp19}%
\begin{bpstatementit}{Corollary~17.6}{ ($H^1_c$ as a stalk of the Fourier transform)}
We have \(H_c^2(\mathbf A^1,j_*M)=0\), and there is a Frobenius-equivariant isomorphism \(H_c^1(\mathbf A^1,j_*M)\simeq(j'_*M')_{0'}\) with the stalk at the origin \(0'\) of the dual line. In particular \(H_c^1(\mathbf A^1,j_*M)\) has weights at most \(w+1\).
\end{bpstatementit}

% proof: proof/corollary-h-1-c-stalk-fourier
\begin{proof}[{Proof}]\phantomsection\label{bp20}%
We compute the stalk at \(0'\) of \(\operatorname{FT}_\psi((j_*M)[1])\) in two ways: by the stalk formula, which gives the cohomology of \(j_*M\), and through the isomorphism with \((j'_*M')[1]\), which gives a stalk of a middle extension.

At the rational origin \(0'\) of the dual line, the kernel \(\mathcal K_\psi\) restricts to the pullback of \(\mathcal L_\psi\) along the constant map to \(0\), which is canonically trivial (on traces this is \(\psi(0)=1\)). The stalk formula of Proposition \bpout{16.5}, which is compatible with Frobenius, therefore gives

\[
\bigl(\operatorname{FT}_\psi((j_*M)[1])\bigr)_{0'}
=R\Gamma_c(\mathbf A^1,j_*M)[2].
\]

We compare this with the stalk of \((j'_*M')[1]\), which is concentrated in degree \(-1\). In degree \(0\) we therefore see that \(H_c^2(\mathbf A^1,j_*M)=0\), and in degree \(-1\) we obtain

\[
H_c^1(\mathbf A^1,j_*M)
\simeq(j'_*M')_{0'}
=(M'_{\bar\eta})^{I_{0'}},
\]

where \(I_{0'}\) is the inertia group at the origin \(0'\) of the dual line acting on the generic fibre of \(M'\), and the last equality is the formula of Lemma \bpout{15.2} for the stalk of a middle extension. By Proposition \bpout{15.5} the weights of this stalk are at most the weight of \(M'\), which is \(w+1\) by Theorem \hyperref[bp16]{17.5}.
\end{proof}


% object: sec-The-affine-line (blueprint/chap-The-Deligne-Laumon-theorem-for-curves/sec-The-affine-line/_section.md), content hidden
\subsection*{\texorpdfstring{17.3\quad The affine line}{17.3 The affine line}}
\phantomsection\label{bp21}%
\addcontentsline{toc}{subsection}{\texorpdfstring{17.3\quad The affine line}{17.3 The affine line}}
% text: prose-17-07
\phantomsection\label{bp22}%
To pass from irreducible sheaves to arbitrary ones we show that \(j_*\) is exact on sequences of pure sheaves of one weight. The reason is a gap in weights. Let \(0\to F_1\to F\to F_2\to0\) be an exact sequence of lisse sheaves on a dense open \(U\) of a curve \(C\). Applying \(j_*\) gives an exact sequence \(0\to j_*F_1\to j_*F\to j_*F_2\to R^1j_*F_1\), and \(R^1j_*F_1\) is supported on \(C\setminus U\). Now suppose that \(F_1\) and \(F_2\) are pure of weight \(w\). Then, by Proposition \bpout{15.5}, the source of the last map has weights at most \(w\) at a point of \(C\setminus U\). On the other hand, by local duality, which we prove first, the stalk of \(R^1j_*F_1\) at such a point is the dual of the invariants of \(F_1^\vee\), twisted by \((-1)\), and this space has weights at least \(w+2\). Since a Frobenius-equivariant map between spaces with disjoint weights is zero, \(j_*\) is exact on such sequences (Proposition \hyperref[bp27]{17.9}).

The hypothesis of equal weights is needed. In the Kummer extension \(0\to\overline{\mathbf Q}_\ell(1)\to F\to\overline{\mathbf Q}_\ell\to0\) of the chapter opening, \(F_1=\overline{\mathbf Q}_\ell(1)\) has weight \(-2\) and \(F_2=\overline{\mathbf Q}_\ell\) has weight \(0\), so at the origin the stalks of \(j_*\overline{\mathbf Q}_\ell\) and of \(R^1j_*\overline{\mathbf Q}_\ell(1)\) both have weight \(0\), and the connecting map is indeed nonzero.

Exactness lets us extend Corollary \hyperref[bp19]{17.6} from one irreducible sheaf with no geometric constituent of the form \(\mathcal L_{\psi(ax)}\) to a pure sheaf unramified at infinity (Theorem \hyperref[bp30]{17.10}) by induction along a composition series. Each factor of the series is either geometrically constant, and then \(H^1_c=0\), or it satisfies the hypotheses of the corollary.

% object: lemma-local-duality-point-c-u (blueprint/chap-The-Deligne-Laumon-theorem-for-curves/sec-The-affine-line/lemma-local-duality-point-c-u/_section.md)
\phantomsection\label{bp23}%
\begin{bpstatementit}{Lemma~17.7}{ (Local duality at a point of $C\setminus U$)}
Let \(j_0:U_0\hookrightarrow C_0\) be a dense open in a smooth curve over \(k_0\), and let \(F\) be a lisse Weil \(\overline{\mathbf Q}_\ell\)-sheaf on \(U_0\). For a geometric point \(\bar s\) over a closed point \(s\in C_0\setminus U_0\), there is a canonical perfect pairing, equivariant for the local geometric Frobenius at \(s\),

\[
(j_*F^\vee)_{\bar s}\otimes(R^1j_*F)_{\bar s}
\longrightarrow\overline{\mathbf Q}_\ell(-1).
\]
\end{bpstatementit}

% proof: proof/lemma-local-duality-point-c-u
\begin{proof}[{Proof}]\phantomsection\label{bp24}%
The two stalks are cohomology sheaves of the complexes \(i^*Rj_*F^\vee\) (in degree zero) and \(i^*Rj_*F\) (in degree one) on the boundary. We compute the Verdier dual of the second complex by duality on the curve, and we find the first one, shifted and twisted.

Let \(i:Z\hookrightarrow C\) be the inclusion of the whole finite complement, and put \(K_F=i^*Rj_*F\). For a lisse sheaf \(G\) on \(U\), the adjunction \(j^*\dashv Rj_*\) shows that \(i^!Rj_*G=0\), and applying \(i^!\) to the localization triangle for \(Rj_*G\) we obtain \(i^!j_!G\simeq i^*Rj_*G[-1]\). Since duality exchanges \(Rj_*\) with \(j_!\) and \(i^*\) with \(i^!\), this formula lets us write the result again in the form \(i^*Rj_*\). Indeed, by Verdier duality on the smooth curve, \(\mathbb D_UF=F^\vee(1)[2]\) for the lisse sheaf \(F\), and we find

\[
\begin{split}
\mathbb D_Z K_F
&\simeq i^!\mathbb D_C Rj_*F
\simeq i^!j_!\mathbb D_UF\\
&\simeq i^*Rj_*(F^\vee(1)[2])[-1]
=K_{F^\vee}(1)[1].
\end{split}
\]

On the finite geometric scheme \(Z\), Verdier duality is the exact duality of coefficient vector spaces, so taking degree \(-1\) at \(\bar s\) we get \((R^1j_*F)_{\bar s}^\vee=(j_*F^\vee)_{\bar s}(1)\), which is the pairing. It is Frobenius-equivariant because the adjunction, localization and duality maps are natural in \(F\) and respect its Weil structure.
\end{proof}

% text: prose-17-08
\phantomsection\label{bp25}%
For \(F\) pure lisse of weight \(w\), the pairing shows that \((R^1j_*F)_{\bar s}\) has weights at least \(w+2\). Indeed, it puts this stalk in duality with a stalk of \(j_*F^\vee\), which has weights at most \(-w\), and the pairing takes values in \(\overline{\mathbf Q}_\ell(-1)\), which has weight \(2\). The next proposition uses this gap between the source and the target of the connecting map.

% object: example-local-duality-origin (blueprint/chap-The-Deligne-Laumon-theorem-for-curves/sec-The-affine-line/example-local-duality-origin/_section.md)
\phantomsection\label{bp26}%
\begin{bpstatement}{Example~17.8}{ (Local duality at the origin of the line)}
Let \(j:\mathbf G_m\hookrightarrow\mathbf A^1\) over \(k_0\), let \(F=\overline{\mathbf Q}_\ell\), which is pure of weight \(0\), and let \(s=0\). Since inertia at the origin acts trivially on the constant sheaf, \((j_*\overline{\mathbf Q}_\ell)_{\bar 0}=\overline{\mathbf Q}_\ell\) by Lemma \bpout{15.2}, with Frobenius acting by \(1\), of weight \(0\). On the other side, the Kummer covers \(t\mapsto t^{\ell^n}\) of a punctured neighbourhood of the origin give \((R^1j_*\overline{\mathbf Q}_\ell)_{\bar 0}=\overline{\mathbf Q}_\ell(-1)\) (Chapter \bpout{6}), on which geometric Frobenius acts by \(q\), of weight \(2\). The pairing of Lemma \hyperref[bp23]{17.7} is the multiplication \(\overline{\mathbf Q}_\ell\otimes\overline{\mathbf Q}_\ell(-1)\to\overline{\mathbf Q}_\ell(-1)\). Thus the two stalks have weights \(0\) and \(2\), and the gap of two comes entirely from the Tate twist. Indeed, with coefficients \(\overline{\mathbf Q}_\ell(1)\), \(H^1\) of a punctured disc is spanned by the Kummer class of the coordinate, which Frobenius fixes, so with coefficients \(\overline{\mathbf Q}_\ell=\overline{\mathbf Q}_\ell(1)\otimes\overline{\mathbf Q}_\ell(-1)\) Frobenius acts by \(q\).
\end{bpstatement}

% object: proposition-exactness-fixed-pure-weight (blueprint/chap-The-Deligne-Laumon-theorem-for-curves/sec-The-affine-line/proposition-exactness-fixed-pure-weight/_section.md)
\phantomsection\label{bp27}%
\begin{bpstatementit}{Proposition~17.9}{ (Exactness at a fixed pure weight)}
Let \(j_0:U_0\hookrightarrow C_0\) be as in Lemma \hyperref[bp23]{17.7}. If \(0\to F_1\to F\to F_2\to0\) is a short exact sequence of lisse Weil sheaves on \(U_0\), with \(F_1\) and \(F_2\) pure of the same weight \(w\), then the sequence

\[
0\longrightarrow j_*F_1\longrightarrow j_*F
\longrightarrow j_*F_2\longrightarrow0
\]

is exact.
\end{bpstatementit}

% proof: proof/proposition-exactness-fixed-pure-weight
\begin{proof}[{Proof}]\phantomsection\label{bp28}%
Applying \(j_*\) to the short exact sequence gives the long exact sequence

\[
0\to j_*F_1\to j_*F\to j_*F_2\xrightarrow{\ \delta\ }R^1j_*F_1,
\]

so it suffices to show that the connecting map \(\delta\) vanishes. It is Frobenius-equivariant, and we show that its source and target have disjoint weights. Since \(R^1j_*F_1\) is zero on \(U\), we look at a point of \(C\setminus U\). There the stalk of the source has weights at most \(w\) by Proposition \bpout{15.5}. For the target, note that the dual \(F_1^\vee\) is pure of weight \(-w\), so the same result shows that the stalk of \(j_*F_1^\vee\) has weights at most \(-w\). By Lemma \hyperref[bp23]{17.7} the stalk of the target is the dual of this stalk twisted by \((-1)\), and hence the target has weights at least \(w+2\). Now the image of \(\delta\) is a quotient of the source and a submodule of the target, so its eigenvalues would have weights both at most \(w\) and at least \(w+2\). Hence \(\delta\) vanishes at each point of \(C\setminus U\). As \(R^1j_*F_1\) is zero on \(U\), it follows that \(\delta=0\).
\end{proof}

% text: prose-17-09
\phantomsection\label{bp29}%
All factors of a composition series of a pure lisse sheaf are pure of the same weight, so the proposition applies to every step of such a series. In this way we pass from the irreducible sheaves of Corollary \hyperref[bp19]{17.6} to arbitrary ones.

% object: theorem-weight-bound-affine-line (blueprint/chap-The-Deligne-Laumon-theorem-for-curves/sec-The-affine-line/theorem-weight-bound-affine-line/_section.md)
\phantomsection\label{bp30}%
\begin{bpstatementit}{Theorem~17.10}{ (The weight bound on the affine line)}
Let \(j:U_0\hookrightarrow\mathbf A^1_{k_0}\) be a dense open, and let \(F\) be a lisse sheaf on \(U_0\), pure of weight \(w\) and unramified at infinity. Then \(H_c^1(\mathbf A^1,j_*F)\) has weights at most \(w+1\).
\end{bpstatementit}

% proof: proof/theorem-weight-bound-affine-line
\begin{proof}[{Proof}]\phantomsection\label{bp31}%
We prove the bound for the irreducible factors of \(F\) and pass to \(F\) by induction. The irreducible factors are of two kinds: geometrically constant ones, for which \(H^1_c\) vanishes, and the others, to which Corollary \hyperref[bp19]{17.6} applies.

Take a finite composition series of the finite-dimensional Weil representation of \(F\). Each irreducible lisse factor is pure of weight \(w\), since its local eigenvalues are eigenvalues of a subquotient. It is again unramified at infinity, since inertia at infinity acts trivially on every subquotient of a representation on which it acts trivially. An irreducible factor that is geometrically constant is a sheaf \(\mathcal B_b\) of Proposition \bpout{13.7}, because the geometric action is trivial and an invertible linear operator over \(\overline{\mathbf Q}_\ell\) has an eigenline. Its middle extension is then constant on the affine line, and so its \(H_c^1\) vanishes by the trivial case of Lemma \bpout{16.4}. Every other factor \(F_i\) is irreducible, unramified at infinity and not geometrically constant, so by Lemma \bpout{16.18} it has no geometric constituent of the form \(\mathcal L_{\psi(ax)}\), and we may apply Corollary \hyperref[bp19]{17.6} with \(M=F_i\).

The induction needs \(j_*\) to turn the composition series into a filtration of \(j_*F\) whose quotients are the \(j_*F_i\). Since all factors are pure of the same weight \(w\), Proposition \hyperref[bp27]{17.9} shows that the short exact sequences of the series remain exact after applying \(j_*\). So consider a step \(0\to F'\to F''\to F_i\to0\) of the series, where we may assume that the bound holds for \(F'\). In the resulting long exact sequence of compactly supported cohomology, the image of \(H^1_c(\mathbf A^1,j_*F')\) in \(H^1_c(\mathbf A^1,j_*F'')\) is a quotient of the former, and the cokernel embeds in \(H^1_c(\mathbf A^1,j_*F_i)\). Since the bound \(w+1\) is preserved under subquotients and extensions, we conclude by induction that it holds for \(F\).
\end{proof}

% text: prose-17-10
\phantomsection\label{bp32}%
The condition at infinity enters only through Lemma \bpout{16.18}, to verify the hypothesis of Proposition \bpout{16.17}: it excludes the constituents \(\mathcal L_{\psi(ax)}\) with \(a\ne0\), whose transforms are skyscrapers rather than middle extensions. In the last section we remove this hypothesis, together with the restriction to opens of the line.


% object: sec-Arbitrary-curves (blueprint/chap-The-Deligne-Laumon-theorem-for-curves/sec-Arbitrary-curves/_section.md), content hidden
\subsection*{\texorpdfstring{17.4\quad Arbitrary curves}{17.4 Arbitrary curves}}
\phantomsection\label{bp33}%
\addcontentsline{toc}{subsection}{\texorpdfstring{17.4\quad Arbitrary curves}{17.4 Arbitrary curves}}
% text: prose-17-11
\phantomsection\label{bp34}%
Theorem \hyperref[bp30]{17.10} assumes that the curve is an open of the affine line and that the sheaf is unramified at infinity. We remove the second hypothesis by moving a point where the sheaf is lisse to infinity, and the first by taking the direct image along a finite map to the line.

For an upper bound it is enough to replace \(H^1_c\) by a group that surjects onto it, and each move does this. For the first move, let \(s\) be a point where the sheaf is lisse. Since a point has no \(H^1_c\), deleting \(s\) gives a surjection \(H^1_c(W\setminus\{s\},F)\to H^1_c(W,F)\), and after a finite extension of the constant field and an automorphism of \(\mathbf P^1\) sending \(s\) to infinity, the sheaf is unramified there. For the second, a finite map \(\pi:C\to\mathbf A^1\) does not change compactly supported cohomology at all, since \(R\Gamma_c(\mathbf A^1,\pi_*j_!F)=R\Gamma_c(C,F)\). However, \(\pi_*j_!F\) is lisse and pure only off the branch locus, so we remove the bad points by one more localization. Such a map \(\pi\) exists by Riemann–Roch, because a function with a single pole of order prime to \(p\) defines a finite generically étale map (Lemma \hyperref[bp41]{17.13}). Curves that are not affine, and the degrees \(r\ne1\), are treated afterwards in Theorem \hyperref[bp47]{17.15}.

% object: lemma-weight-bound-constant-extension (blueprint/chap-The-Deligne-Laumon-theorem-for-curves/sec-Arbitrary-curves/lemma-weight-bound-constant-extension/_section.md)
\phantomsection\label{bp35}%
\begin{bpstatementit}{Lemma~17.11}{ (Weight bounds and finite constant extensions)}
Let \(X_0\) be a separated scheme of finite type over \(k_0\), let \(F\) be a constructible Weil sheaf on \(X_0\), let \(k_m\supset k_0\) be the extension of degree \(m\), and let \(u\in\mathbf R\). The group \(H^r_c(X,F)\) is the same for \(X_0\) and for \(X_0\otimes_{k_0}k_m\), and the geometric Frobenius of \(k_m\) acts on it as the \(m\)th power of that of \(k_0\). Every eigenvalue \(\alpha\) of the Frobenius of \(k_0\) on \(H^r_c(X,F)\) satisfies \(|\iota\alpha|\le q^{u/2}\) if and only if every eigenvalue \(\beta\) of the Frobenius of \(k_m\) satisfies \(|\iota\beta|\le q^{mu/2}\).
\end{bpstatementit}

% proof: proof/lemma-weight-bound-constant-extension
\begin{proof}[{Proof}]\phantomsection\label{bp36}%
An extension of degree \(m\) of the constant field does not change the geometric scheme \(X\), the sheaf on it, or its cohomology, and it replaces the Frobenius of \(k_0\) by its \(m\)th power. So the eigenvalues \(\beta\) are the \(\alpha^m\), and \(|\iota\alpha^m|\le q^{mu/2}\) holds if and only if \(|\iota\alpha|\le q^{u/2}\).
\end{proof}

% text: prose-17-12
\phantomsection\label{bp37}%
To move a point to infinity we need a rational point, and hence a finite extension of the constant field. By the lemma this does not affect a weight bound.

% object: lemma-finite-extension-base-field-moving (blueprint/chap-The-Deligne-Laumon-theorem-for-curves/sec-Arbitrary-curves/lemma-finite-extension-base-field-moving/_section.md)
\phantomsection\label{bp38}%
\begin{bpstatementit}{Lemma~17.12}{ (Moving a lisse point to infinity)}
Let \(W_0\) be a dense open of \(\mathbf P^1_{k_0}\), and let \(F\) be a lisse Weil sheaf on \(W_0\), pure of weight \(w\). There are a finite extension \(k_m\) of \(k_0\), a dense open \(V_0\) of \(\mathbf A^1_{k_m}\), a lisse Weil sheaf \(G\) on \(V_0\), pure of weight \(w\) and unramified at infinity, and a surjection

\[
H_c^1(V,G)\twoheadrightarrow H_c^1(W,F)
\]

that commutes with the geometric Frobenius of \(k_m\).
\end{bpstatementit}

% proof: proof/lemma-finite-extension-base-field-moving
\begin{proof}[{Proof}]\phantomsection\label{bp39}%
We delete a rational point where \(F\) is lisse, which replaces \(H^1_c\) by a group surjecting onto it, and move that point to infinity by an automorphism of \(\mathbf P^1\). First choose a closed point of \(W_0\) and a finite extension \(k_m\) of the constant field over which it has a rational point \(s\). The restriction of \(F\) to \(W_0\otimes_{k_0}k_m\) is again pure of weight \(w\), because a residue extension raises an eigenvalue and the residue cardinality to the same power. Next, some element of \(\mathrm{PGL}_2(k_m)\) sends \(s\) to infinity, and we transport \(F|_{W\setminus\{s\}}\) by this automorphism. As \(F\) is lisse at \(s\), its local inertia there is trivial, so the transported sheaf \(G\) extends to a lisse sheaf across the new infinity. It is defined on the image \(V_0\) of \(W_0\setminus\{s\}\), a dense open of the affine line, and \(H^1_c(V,G)=H^1_c(W\setminus\{s\},F)\) because the automorphism identifies the two pairs. The surjection follows from the localization sequence

\[
H_c^1(W\setminus\{s\},F)\to H_c^1(W,F)\to H_c^1(\{s\},F_{\bar s})=0,
\]

in which the last group vanishes because a point has no cohomology in positive degrees. Since every construction is defined over \(k_m\), the map commutes with the Frobenius of \(k_m\).
\end{proof}

% text: prose-17-13
\phantomsection\label{bp40}%
By the lemma, for any open of the projective line and any pure lisse sheaf on it, we may bound \(H^1_c\) with Theorem \hyperref[bp30]{17.10}. To reach an arbitrary affine curve we map it to the line, and we have to choose the map so that the direct image is again pure and lisse on a dense open.

% object: lemma-finite-generically-etale-map-line (blueprint/chap-The-Deligne-Laumon-theorem-for-curves/sec-Arbitrary-curves/lemma-finite-generically-etale-map-line/_section.md)
\phantomsection\label{bp41}%
\begin{bpstatementit}{Lemma~17.13}{ (A finite generically étale map to the line)}
Let \(C_0\) be a smooth geometrically connected affine curve. There are a smooth curve \(C'_0\) containing \(C_0\) with finite complement and a finite generically étale map \(\pi:C'_0\to\mathbf A^1_{k_0}\). If \(F\) is a pure lisse sheaf of weight \(w\) on \(C_0\), then the sheaf \(K=\pi_*j_!F\), where \(j:C_0\hookrightarrow C'_0\), is lisse and pure of weight \(w\) outside the branch locus and the image of \(C'_0\setminus C_0\). Moreover, \(R\Gamma_c(\mathbf A^1,K)=R\Gamma_c(C,F)\) canonically.
\end{bpstatementit}

% proof: proof/lemma-finite-generically-etale-map-line
\begin{proof}[{Proof}]\phantomsection\label{bp42}%
The map is given by a function with a single pole, at a point at infinity, and we choose the order of this pole prime to \(p\) so that the map is generically étale. We find the function by Riemann–Roch, we show by a local computation at the pole that its differential is nonzero, and we check the assertions about \(K\) stalk by stalk.

Let \(\overline C_0\) be the smooth projective completion, choose a closed point \(s\) of \(\overline C_0\) outside \(C_0\), and choose \(m\) prime to \(p\) with \((m-1)\deg(s)>2g(\overline C_0)-2\). Both divisors \(ms\) and \((m-1)s\) then have degree \(>2g-2\), so Riemann–Roch (Theorem \bpout{C.3}) gives \(h^0(\mathcal O(ns))=n\deg(s)+1-g\) for \(n=m,m-1\), and hence

\[
h^0(\mathcal O(ms))-h^0(\mathcal O((m-1)s))=\deg(s)>0.
\]

Therefore there is a function \(f\) whose only pole is at \(s\), of exact order \(m\). The sections \(1,f\) define a nonconstant map \(\overline C_0\to\mathbf P^1\) whose inverse image of infinity is \(\{s\}\), and this map is proper with finite fibres, hence finite. To see that it is generically étale, choose a uniformizer \(t\) at \(s\) and write \(f=t^{-m}u\) with \(u\) a unit. Since \(\kappa(s)/k_0\) is separable and the curve is smooth, \(dt\) generates the module of differentials near \(s\). Writing \(du=u'\,dt\), we get

\[
df=t^{-m-1}(-mu+tu')\,dt\ne0,
\]

since \(-mu+tu'\) is congruent to the unit \(-mu\) modulo \(t\), and \(mu\) is a unit because \(p\nmid m\). The differential of this finite map between smooth curves is therefore invertible at the generic point, and the map is generically étale. We set \(C'_0=\overline C_0\setminus\{s\}\) and let \(\pi\) be the restriction of the map to \(C'_0\).

Over the open in the statement, \(\pi\) is finite étale and its inverse image lies in \(C_0\), where \(j_!F=F\). Moreover, the direct image of a lisse sheaf under a finite étale map is lisse, because étale locally the map is a disjoint union of copies of the base. Hence \(K\) is lisse there. For purity, note that at a closed point \(y\) of this open, the stalk of \(K\) is a sum over the geometric points of the fibre, which Frobenius permutes. Consider such a cycle of length \(e\), corresponding to a point over \(y\) with residue extension of degree \(e\). On it the \(e\)th power of the operator is the local Frobenius upstairs, so an eigenvalue \(\beta\) satisfies \(|\iota\beta|^e=q_y^{ew/2}\), and the weight is still \(w\). Since \(R\pi_!=\pi_*\) for the finite map \(\pi\), the identification of cohomology follows by composing direct images with compact support. It commutes with Frobenius because the construction is defined over \(k_0\).
\end{proof}

% text: prose-17-14
\phantomsection\label{bp43}%
For example, take for \(C_0\) the affine curve \(y^2=x^3-x\) over \(\mathbf F_3\), the complement of the point at infinity of the projective curve. There \(g=1\), the point \(s=\infty\) is rational, and \(m=2\) satisfies \((m-1)\deg(s)>0\). Accordingly, the function \(x\) has a single pole, of order two, which is prime to \(3\), and the map is the double cover \((x,y)\mapsto x\) of the line, whose extension to the projective curve is branched over \(0\), \(\pm1\) and \(\infty\).

With both lemmas the bound on the line extends to every affine curve.

% object: theorem-weight-bound-affine-curve (blueprint/chap-The-Deligne-Laumon-theorem-for-curves/sec-Arbitrary-curves/theorem-weight-bound-affine-curve/_section.md)
\phantomsection\label{bp44}%
\begin{bpstatementit}{Theorem~17.14}{ (The weight bound on an affine curve)}
For a pure lisse Weil sheaf \(F\) of weight \(w\) on a smooth geometrically connected affine curve \(C_0\), the group \(H_c^1(C,F)\) has weights at most \(w+1\).
\end{bpstatementit}

% proof: proof/theorem-weight-bound-affine-curve
\begin{proof}[{Proof}]\phantomsection\label{bp45}%
We push the sheaf forward to the line and realize \(H^1_c\) as a quotient of a group bounded by Theorem \hyperref[bp30]{17.10}. This uses two surjections: one removes the points where the direct image fails to be lisse and pure, and the other moves a lisse point to infinity.

Choose \(\pi\) as in Lemma \hyperref[bp41]{17.13} and put \(K=\pi_*j_!F\). Let \(W_0\) be the dense open on which \(K\) is lisse and pure. Since its complement is finite, localization gives

\[
H_c^1(W,K|_W)\twoheadrightarrow
H_c^1(\mathbf A^1,K)=H_c^1(C,F).
\]

By Lemma \hyperref[bp38]{17.12}, after a finite extension of the constant field there is a further surjection from \(H_c^1(V,G)\) onto \(H_c^1(W,K|_W)\), with \(G\) pure of weight \(w\) and unramified at infinity. Let \(v:V\hookrightarrow\mathbf A^1\) be the inclusion and \(i:Z\hookrightarrow\mathbf A^1\) its finite complement. The localization sequence for \(v_*G\) reads

\[
H_c^0(Z,i^*v_*G)\longrightarrow H_c^1(V,G)
\longrightarrow H_c^1(\mathbf A^1,v_*G)\longrightarrow0.
\]

The first group has weights at most \(w\). Indeed, the stalks of \(i^*v_*G\) at the closed points of \(Z_0\) have weights at most \(w\) by Proposition \bpout{15.5}, and \(H_c^0(Z,i^*v_*G)\) is the direct sum, over the closed points \(z\) of \(Z_0\), of the stalks at the \(e=[\kappa(z):k_m]\) geometric points over \(z\), where \(k_m\) is the constant field of \(V_0\). The geometric Frobenius of \(k_m\) permutes these \(e\) stalks cyclically, and its \(e\)th power on each of them is conjugate to the local Frobenius \(F_z\). An eigenvalue \(\beta\) on such a cycle therefore satisfies \(\beta^e=\alpha\) for an eigenvalue \(\alpha\) of \(F_z\), so that \(|\iota\beta|^e=|\iota\alpha|\le q_z^{w/2}=(q^m)^{ew/2}\), and \(|\iota\beta|\le(q^m)^{w/2}\), which is weight at most \(w\) over \(k_m\). This is the case of the finite map \(Z_0\to\operatorname{Spec}k_m\) in Lemma \bpout{18.1}, whose proof is the same computation. The last group has weights at most \(w+1\) by Theorem \hyperref[bp30]{17.10}. As an extension of a subgroup of the last by a quotient of the first, the middle group therefore has weights at most \(w+1\). This bound passes to quotients under the two surjections, and by Lemma \hyperref[bp35]{17.11} it holds over the original constant field.
\end{proof}

% text: prose-17-15
\phantomsection\label{bp46}%
To reach a curve that is not affine, we pass to a dense affine open and view the sheaf on the whole curve as a middle extension.

% object: theorem-weight-bound-curve (blueprint/chap-The-Deligne-Laumon-theorem-for-curves/sec-Arbitrary-curves/theorem-weight-bound-curve/_section.md)
\phantomsection\label{bp47}%
\begin{bpstatementit}{Theorem~17.15}{ (The weight bound on a curve)}
Let \(j:U_0\hookrightarrow C_0\) be a dense open in a smooth curve over \(k_0\), and let \(F\) be a pure lisse Weil sheaf of weight \(w\) on \(U_0\). Then for every \(r\in\mathbf Z\) the group \(H_c^r(C,j_*F)\) has weights at most \(w+r\). In particular, for \(U_0=C_0\), the group \(H_c^r(C,F)\) has weights at most \(w+r\).
\end{bpstatementit}

% proof: proof/theorem-weight-bound-curve
\begin{proof}[{Proof}]\phantomsection\label{bp48}%
On a geometrically connected curve we compare \(j_*F\) with its restriction to a dense affine open, where the affine case applies, and we deduce the bounds from the localization sequence. The general case reduces to this one, because a curve that is not geometrically connected becomes a disjoint union of geometrically connected ones after a finite constant extension.

Suppose first that \(C_0\) is geometrically connected, and choose a dense affine open \(v:V_0\hookrightarrow U_0\). Since \(F\) is lisse on \(U_0\), the unit \(F\to v_*F|_V\) is an isomorphism by the last assertion of Lemma \bpout{15.2}. Hence \(j_*F\) is also the middle extension of \(F|_V\) from \(V_0\) to \(C_0\). On \(V\), by Theorem \hyperref[bp44]{17.14} and Lemma \bpout{15.7}, \(H_c^r(V,F)\) has weights at most \(w+r\) for every \(r\), and \(H_c^0(V,F)=0\). On the complement, write \(Z_0=C_0\setminus V_0\) and \(i\) for its inclusion; then the restriction of \(j_*F\) to the finite set \(Z_0\) has weights at most \(w\) by Proposition \bpout{15.5}, and its cohomology is concentrated in degree zero. These two pieces are joined by the localization sequence for \(V\subset C\), which reads

\[
0=H_c^0(V,F)\to H_c^0(C,j_*F)\to H^0(Z,i^*j_*F)\to H_c^1(V,F)
\to H_c^1(C,j_*F)\to0,
\]

followed by an isomorphism \(H_c^2(V,F)\simeq H_c^2(C,j_*F)\), and all higher groups vanish. Over \(\operatorname{Spec}k_0\) the bounds are conditions on eigenvalues (Lemmas \bpout{13.2} and \bpout{Lemma 13.3}), and we read them off degree by degree. In degree zero, the group \(H_c^0(C,j_*F)\) injects into \(H^0(Z,i^*j_*F)\), which has weights at most \(w\): Frobenius permutes the stalks at the geometric points over each closed point of \(Z_0\) cyclically, with the local Frobenius as the power that returns to a stalk, as in the proof of Theorem \hyperref[bp44]{17.14}. In degree one, \(H_c^1(C,j_*F)\) is a quotient of \(H_c^1(V,F)\), which has weights at most \(w+1\). In degree two the groups are isomorphic, and the bound \(w+2\) holds for \(V\).

In general, the connected components of \(C\) are defined over a finite extension \(k_m\) of \(k_0\), and over \(k_m\) each of them is a smooth geometrically connected curve. For example, \(\mathbf P^1_{\mathbf F_{q^2}}\), viewed as a curve over \(\mathbf F_q\), is connected, but \(C\) has two components, which Frobenius exchanges and which are defined over \(\mathbf F_{q^2}\). As Frobenius permutes the components, only its \(m\)th power acts on the cohomology of each of them. On such a component, \(j_*F\) is the middle extension of the restriction of \(F\) to the dense open where the component meets \(U\), and this restriction is pure of weight \(w\). Moreover, the group \(H^r_c(C,j_*F)\) is the direct sum of the groups of these components, compatibly with the \(m\)th power of Frobenius. Hence the bound holds over \(k_m\) by the first case, and by Lemma \hyperref[bp35]{17.11} it holds over \(k_0\).
\end{proof}



\clearpage\phantomsection\hypertarget{bp-outside}{}%
\section*{Outside this export}
\addcontentsline{toc}{section}{Outside this export}
References marked \ensuremath{{}^\ast} are to parts of the blueprint this document leaves out:
\begin{description}
\item[Chapter 6] Constructibility of $f^!$, in Ⅰ The six operations (\bpslug{chap-{}\allowbreak{}Constructibility-{}\allowbreak{}of-{}\allowbreak{}f})
\item[Corollary 12.26] Trace functions and direct image with compact support, in Ⅱ Duality and the trace formula (\bpslug{corollary-{}\allowbreak{}trace-{}\allowbreak{}functions-{}\allowbreak{}direct-{}\allowbreak{}image-{}\allowbreak{}compact})
\item[Chapter 13] Weights, in Ⅲ Weights and the Riemann hypothesis (\bpslug{chap-{}\allowbreak{}Weights})
\item[Lemma 13.2] Weil sheaves on a point are mixed, in Ⅲ Weights and the Riemann hypothesis (\bpslug{lemma-{}\allowbreak{}weil-{}\allowbreak{}sheaves-{}\allowbreak{}point-{}\allowbreak{}mixed})
\item[Lemma 13.3] Bounding the weights of a mixed sheaf, in Ⅲ Weights and the Riemann hypothesis (\bpslug{lemma-{}\allowbreak{}bounding-{}\allowbreak{}weights-{}\allowbreak{}mixed-{}\allowbreak{}sheaf})
\item[Proposition 13.7] Weights under the usual operations, in Ⅲ Weights and the Riemann hypothesis (\bpslug{proposition-{}\allowbreak{}weights-{}\allowbreak{}under-{}\allowbreak{}usual-{}\allowbreak{}operations})
\item[Theorem 14.26] Purity of real sheaves on curves, in Ⅲ Weights and the Riemann hypothesis (\bpslug{theorem-{}\allowbreak{}purity-{}\allowbreak{}real-{}\allowbreak{}sheaves-{}\allowbreak{}curves})
\item[Chapter 15] Pure sheaves on curves, in Ⅲ Weights and the Riemann hypothesis (\bpslug{chap-{}\allowbreak{}Pure-{}\allowbreak{}sheaves-{}\allowbreak{}on-{}\allowbreak{}curves})
\item[Lemma 15.2] The boundary stalk, in Ⅲ Weights and the Riemann hypothesis (\bpslug{lemma-{}\allowbreak{}boundary-{}\allowbreak{}stalk})
\item[Proposition 15.5] Weights of $j_*$ at the boundary points, in Ⅲ Weights and the Riemann hypothesis (\bpslug{proposition-{}\allowbreak{}weights-{}\allowbreak{}j-{}\allowbreak{}boundary-{}\allowbreak{}points})
\item[Corollary 15.6] An unramified boundary point, in Ⅲ Weights and the Riemann hypothesis (\bpslug{corollary-{}\allowbreak{}unramified-{}\allowbreak{}boundary-{}\allowbreak{}point})
\item[Lemma 15.7] The extreme degrees on an affine curve, in Ⅲ Weights and the Riemann hypothesis (\bpslug{lemma-{}\allowbreak{}extreme-{}\allowbreak{}degrees-{}\allowbreak{}affine-{}\allowbreak{}curve})
\item[Proposition 15.9] Twisting to weight zero and adding the dual, in Ⅲ Weights and the Riemann hypothesis (\bpslug{proposition-{}\allowbreak{}twisting-{}\allowbreak{}weight-{}\allowbreak{}zero-{}\allowbreak{}adding-{}\allowbreak{}dual})
\item[Definition 15.10] Growth of squared traces, in Ⅲ Weights and the Riemann hypothesis (\bpslug{definition-{}\allowbreak{}growth-{}\allowbreak{}squared-{}\allowbreak{}traces})
\item[Corollary 15.15] Sums of squared traces of $j_!$ and $j_*$, in Ⅲ Weights and the Riemann hypothesis (\bpslug{corollary-{}\allowbreak{}sums-{}\allowbreak{}squared-{}\allowbreak{}traces-{}\allowbreak{}j-{}\allowbreak{}shriek})
\item[Chapter 16] The Fourier transform, in Ⅲ Weights and the Riemann hypothesis (\bpslug{chap-{}\allowbreak{}The-{}\allowbreak{}Fourier-{}\allowbreak{}transform})
\item[Lemma 16.4] The integral of a linear character, in Ⅲ Weights and the Riemann hypothesis (\bpslug{lemma-{}\allowbreak{}integral-{}\allowbreak{}linear-{}\allowbreak{}character})
\item[Proposition 16.5] Base change and stalks, in Ⅲ Weights and the Riemann hypothesis (\bpslug{proposition-{}\allowbreak{}base-{}\allowbreak{}change-{}\allowbreak{}fourier-{}\allowbreak{}inversion})
\item[Corollary 16.8] Plancherel, in Ⅲ Weights and the Riemann hypothesis (\bpslug{corollary-{}\allowbreak{}plancherel})
\item[Lemma 16.14] The simple objects, in Ⅲ Weights and the Riemann hypothesis (\bpslug{lemma-{}\allowbreak{}simple-{}\allowbreak{}objects})
\item[Proposition 16.17] The Fourier transform of a middle extension, in Ⅲ Weights and the Riemann hypothesis (\bpslug{proposition-{}\allowbreak{}fourier-{}\allowbreak{}transform-{}\allowbreak{}middle-{}\allowbreak{}extension})
\item[Lemma 16.18] Sheaves unramified at infinity, in Ⅲ Weights and the Riemann hypothesis (\bpslug{lemma-{}\allowbreak{}unramified-{}\allowbreak{}at-{}\allowbreak{}infinity-{}\allowbreak{}no-{}\allowbreak{}linear-{}\allowbreak{}characters})
\item[Lemma 18.1] Finite and quasi-finite direct images, in Ⅲ Weights and the Riemann hypothesis (\bpslug{lemma-{}\allowbreak{}finite-{}\allowbreak{}quasi-{}\allowbreak{}finite-{}\allowbreak{}direct-{}\allowbreak{}images})
\item[Theorem C.3] Riemann–Roch and coherent Serre duality, in Appendices (\bpslug{theorem-{}\allowbreak{}riemann-{}\allowbreak{}roch-{}\allowbreak{}coherent-{}\allowbreak{}serre-{}\allowbreak{}duality})
\end{description}
\clearpage\phantomsection%
\section*{References}
\addcontentsline{toc}{section}{References}
\bpbibitem{BBD82}{bp49}{A. Beilinson, J. Bernstein, and P. Deligne. \emph{Faisceaux pervers}, volume 100 of \emph{Astérisque}. 1982.}

\bpbibitem{BS15}{bp50}{B. Bhatt and P. Scholze. The pro-étale topology for schemes. \emph{Astérisque}, 369:99–201, 2015. \href{https://people.mpim-bonn.mpg.de/scholze/proetale.pdf}{Author's manuscript}.}

\bpbibitem{Con07}{bp51}{B. Conrad. Deligne's notes on Nagata compactifications. \emph{Journal of the Ramanujan Mathematical Society}, 22:205–257, 2007. \href{https://math.stanford.edu/\~conrad/papers/nagatafinal.pdf}{Author's manuscript}.}

\bpbibitem{Del74}{bp52}{P. Deligne. La conjecture de Weil. I. \emph{Publ. Math. IHÉS}, 43:273–307, 1974.}

\bpbibitem{Del80}{bp53}{P. Deligne. La conjecture de Weil. II. \emph{Publ. Math. IHÉS}, 52:137–252, 1980.}

\bpbibitem{Den22}{bp54}{C. Deninger. There is no “Weil-”cohomology theory with real coefficients for arithmetic curves. \emph{Ann. Sc. Norm. Super. Pisa Cl. Sci.}, 2022. To appear. arXiv:\href{https://arxiv.org/abs/2204.02714}{2204.02714}.}

\bpbibitem{Deu41}{bp55}{M. Deuring. Die Typen der Multiplikatorenringe elliptischer Funktionenkörper. \emph{Abh. Math. Sem. Univ. Hamburg}, 14:197–272, 1941.}

\bpbibitem{Eke90}{bp56}{T. Ekedahl. On the adic formalism. In \emph{The Grothendieck Festschrift, Vol. II}, volume 87 of \emph{Progress in Mathematics}, pages 197–218. Birkhäuser, 1990.}

\bpbibitem{Has36}{bp57}{H. Hasse. Zur Theorie der abstrakten elliptischen Funktionenkörper. I, II, III. \emph{J. Reine Angew. Math.}, 175:55–62, 69–88, 193–208, 1936.}

\bpbibitem{Hum72}{bp58}{J. E. Humphreys. \emph{Introduction to Lie Algebras and Representation Theory}, volume 9 of \emph{Graduate Texts in Mathematics}. Springer, 1972.}

\bpbibitem{Ill94}{bp59}{L. Illusie. Autour du théorème de monodromie locale. \emph{Astérisque}, 223:9–57, 1994. §4. \href{https://www.numdam.org/item/AST_1994__223__9_0/}{Original article}.}

\bpbibitem{Kat88}{bp60}{N. M. Katz. Travaux de Laumon. \emph{Astérisque}, 161–162:105–132, 1988. Séminaire Bourbaki, exp. 691.}

\bpbibitem{KW01}{bp61}{R. Kiehl and R. Weissauer. \emph{Weil conjectures, perverse sheaves and \(\ell\)-adic Fourier transform}, volume 42 of \emph{Ergebnisse der Mathematik}. Springer, 2001.}

\bpbibitem{Lau87}{bp62}{G. Laumon. Transformation de Fourier, constantes d'équations fonctionnelles et conjecture de Weil. \emph{Publ. Math. IHÉS}, 65:131–210, 1987.}

\bpbibitem{Mil86a}{bp63}{J. S. Milne. \href{https://www.jmilne.org/math/xnotes/AVs.pdf}{Abelian Varieties}. In \emph{Arithmetic Geometry}. 1986. Corrected online version.}

\bpbibitem{Mil86b}{bp64}{J. S. Milne. \href{https://www.jmilne.org/math/xnotes/JVs.pdf}{Jacobian Varieties}. In \emph{Arithmetic Geometry}. 1986. Corrected online version.}

\bpbibitem{Ran39}{bp65}{R. A. Rankin. Contributions to the theory of Ramanujan's function \(\tau(n)\) and similar arithmetical functions. II. \emph{Proc. Cambridge Philos. Soc.}, 35:357–372, 1939.}

\bpbibitem{SGA4}{bp66}{M. Artin, A. Grothendieck, and J.-L. Verdier, editors. \emph{Théorie des topos et cohomologie étale des schémas}, volumes 269, 270, 305 of \emph{Lecture Notes in Mathematics}. Springer, 1972–1973. SGA 4.}

\bpbibitem{SGA4½}{bp67}{P. Deligne. \emph{Cohomologie étale}, volume 569 of \emph{Lecture Notes in Mathematics}. Springer, 1977. SGA 4½; \emph{Théorèmes de finitude}, §3, and \emph{Rapport sur la formule des traces}, §3, and \emph{Sommes trigonométriques}, §1.}

\bpbibitem{SGA5}{bp68}{A. Grothendieck, editor. \emph{Cohomologie \(\ell\)-adique et fonctions \(L\)}, volume 589 of \emph{Lecture Notes in Mathematics}. Springer, 1977. SGA 5.}

\bpbibitem{SGA7II}{bp69}{P. Deligne and N. Katz, editors. \emph{Groupes de monodromie en géométrie algébrique}, volume 340 of \emph{Lecture Notes in Mathematics}. Springer, 1973. SGA 7 II; Exposé XIII, (2.1.7).}

\bpbibitem{Stacks}{bp70}{The Stacks Project Authors. \href{https://stacks.math.columbia.edu}{The Stacks Project}. Individual results are cited by their stable tags in the text.}

\bpbibitem{Var05}{bp71}{Y. Varshavsky. Lefschetz–Verdier trace formula and a generalization of a theorem of Fujiwara. 2005. Version 2. arXiv:\href{https://arxiv.org/abs/math/0505564}{math/0505564}. \href{https://arxiv.org/abs/math/0505564v2}{Author's preprint}.}

\bpbibitem{Wei48a}{bp72}{A. Weil. \emph{Sur les courbes algébriques et les variétés qui s'en déduisent}. Hermann, Paris, 1948.}

\bpbibitem{Wei48b}{bp73}{A. Weil. \emph{Variétés abéliennes et courbes algébriques}. Hermann, Paris, 1948.}

\bpbibitem{Wei49}{bp74}{A. Weil. Numbers of solutions of equations in finite fields. \emph{Bull. Amer. Math. Soc.}, 55:497–508, 1949.}

\end{document}
